\chapter{Le cadre prudentiel de capital et son articulation avec DORA}
\label{chap:reglementaire}

% HARNAIS-HORS-SCRIPT: cadre reglementaire : numeros d'articles, delais de notification et seuils cites des textes (DORA, Solvabilite II, Bale). Ce sont des donnees de droit, pas des resultats.

% MIGRÉ depuis memoire/main.tex (plages 349-617).
% Hiérarchie descendue d'un cran, labels préfixés "reg:".
% RELECTURE FAITE (août 2026). Un résidu de l'ancien modèle a été trouvé et
% corrigé : l'introduction annonçait la non-conformité « comme une variable
% latente », formulation de l'ancien modèle à briques, remplacée par l'état de
% conformité par pilier du ch. 11. Aucune autre phrase du chapitre n'annonce la
% copule ni la latente comme LE modèle.
%==============================================================

Le risque modélisé ici, risque opérationnel d'origine technologique et manquement aux
exigences de résilience numérique, se situe à l'intersection de trois cadres construits
successivement selon des philosophies distinctes : le dispositif bancaire de Bâle, la directive
Solvabilité~II propre à l'assurance, et le règlement transsectoriel DORA. Leur articulation
ouvre un \emph{pont quantitatif} entre exigence de capital et résilience opérationnelle, et
permet de poser la non-conformité comme un état par domaine de contrôle, gradué et évolutif,
dont le chapitre~\ref{chap:multietats} tire la dynamique et le chapitre~\ref{chap:resultats}
le coût en capital.

%--------------------------------------------------------------

%--------------------------------------------------------------

\paragraph{Le risque opérationnel : une définition partagée, des traitements divergents.}
Le risque opérationnel reçoit la même définition dans les secteurs bancaire et
assurantiel. Défini par le Comité de Bâle puis repris par la directive Solvabilité~II,
il désigne le risque de perte résultant de l'inadaptation ou de la défaillance de
procédures internes, de personnels et de systèmes, ou d'événements extérieurs. Cette
définition inclut le risque juridique mais exclut les risques stratégiques et de
réputation.

Si la définition est partagée, les modes de quantification ont divergé. Le secteur bancaire a
abandonné les modèles internes au profit d'une approche standardisée, tandis que l'assurance a
conservé une liberté méthodologique. Cette asymétrie commande le choix méthodologique du
mémoire : Solvabilité~II laisse subsister un espace de modélisation interne, où la
quantification du risque cyber par une distribution de pertes reste recevable.

\paragraph{Bâle III/IV.} Le secteur bancaire a retiré en $2022$ les modèles internes de mesure
avancée au profit d'une approche standardisée, sans invalider l'approche par distribution des
pertes comme outil de modélisation : le mémoire en reprend les deux marges, fréquence et
sévérité, et en remplace l'agrégation par une cascade dirigée. Le détail est déporté
\matcomp.

\paragraph{Solvabilité II : l'asymétrie réglementaire et le maintien du choix méthodologique.}
Contrairement au secteur bancaire, le secteur de l'assurance conserve sous Solvabilité~II
une liberté de choix pour la quantification du risque opérationnel : formule standard,
paramètres propres à l'entreprise (USP), ou modèle interne, partiel ou global, approuvé
par l'autorité de contrôle. Dans la formule standard, l'exigence pour risque opérationnel
($SCR_{Op}$) s'ajoute au capital de solvabilité requis de base
($BSCR$), ajusté de la capacité d'absorption des pertes :
\begin{equation}
SCR = BSCR + Adj + SCR_{Op},
\qquad
SCR_{Op} = \min(0{,}3 \times BSCR,\ Op) + 0{,}25 \times Exp_{ul},
\label{reg:eq:scrop}
\end{equation}
où $Op = \max(Op_{premiums}, Op_{provisions})$ est calculé à partir de coefficients
forfaitaires appliqués aux primes et provisions, vie et non-vie. La forme donnée ici est
simplifiée ; la formule complète de l'Article~204 du règlement délégué~2015/35
comporte en outre des termes de croissance, sous la forme
$\max\!\left(0;\, 0{,}04 \times \Delta\!\right)$, pénalisant les hausses de primes
supérieures à $20\,\%$ d'un exercice à l'autre.

Fondée sur les seuls indicateurs de volume, cette approche pose que deux assureurs gérant le
même volume de primes ont la même exposition au risque opérationnel, quelles que soient la
qualité de leur infrastructure informatique, l'automatisation de leurs processus ou leur
dépendance à des prestataires d'externalisation critiques. C'est cette dépendance
technologique que DORA encadre, et dont le mémoire mesure l'effet sur le capital.

\paragraph{Pilier 1 et Pilier 2.} Le $SCR_{Op}$ relève du Pilier~1, quantitatif et
forfaitaire. Mais Solvabilité~II comporte un second pilier, l'évaluation interne des
risques et de la solvabilité (\emph{Own Risk and Solvency Assessment}, ORSA), qui impose
à l'assureur une évaluation prospective de l'adéquation entre son profil de risque réel
et ses besoins de solvabilité. C'est dans le cadre de l'ORSA, et non du forfait de
Pilier~1 plafonné à $30\,\%$ du $BSCR$, que se situe la modélisation du
mémoire : l'ORSA admet, et encourage, les scénarios de stress stochastiques que la formule
standard ne sait pas produire.

L'additivité que l'équation ci-dessus donne à la charge opérationnelle n'est pas un détail
de présentation. Elle décide de ce que le capital total retient d'un besoin de nature
opérationnelle, et elle le fait dans un sens qui n'est pas neutre pour l'entité : le
\S\ref{sec:agregation-scr} le chiffre sur les quatre bilans réels, et montre que le plafond
lui-même borne ce qu'un forfait saurait exprimer.

\paragraph{DORA, rappelé en une ligne et non réexposé.} Le règlement sur la résilience
opérationnelle numérique du secteur financier, ses cinq domaines de maîtrise et son dispositif de
surveillance des prestataires tiers critiques sont présentés au \S\ref{sec:dora}. Ce chapitre ne
les redit pas : il traite de ce que le chapitre~\ref{chap:cadre-cyber-dora} laisse ouvert, à
savoir l'articulation du règlement avec les exigences de capital qui, elles, préexistent.

\paragraph{Position dans l'écosystème.} La progression Bâle~$\rightarrow$~Solvabilité~II
$\rightarrow$~DORA traduit un déplacement du centre de gravité prudentiel : de la
\emph{couverture financière} du risque (détenir assez de capital) vers la \emph{continuité
opérationnelle} (rester capable d'opérer). DORA ne remplace pas les exigences de capital ;
il en révèle l'angle mort, à savoir la dynamique propre des défaillances technologiques,
que ni le $SCR_{Op}$ forfaitaire ni l'ORSA qualitatif ne savaient capturer.

%--------------------------------------------------------------

%--------------------------------------------------------------

\paragraph{ORSA et DORA.} Les deux cadres partagent un objet, le risque opérationnel d'origine
technologique, et une exigence d'évaluation prospective, mais pas leur finalité : Solvabilité~II
vise la continuité financière, DORA la continuité opérationnelle. La contribution du mémoire se
loge dans leur recouvrement, en utilisant les artefacts produits par DORA pour quantifier dans
l'ORSA ce que celui-ci ne savait que décrire. Les limites de l'ORSA traditionnel, la table
comparative des deux cadres et leurs zones de divergence sont déportées \matcomp.

%--------------------------------------------------------------

% =====================================================================
